%\documentclass[draft]{article}
\documentclass{article}

\usepackage{arxiv}

\usepackage[utf8]{inputenc} % allow utf-8 input
\usepackage[T1]{fontenc}    % use 8-bit T1 fonts
\usepackage{hyperref}       % hyperlinks
\usepackage{url}            % simple URL typesetting
\usepackage{booktabs}       % professional-quality tables
\usepackage{amsfonts}       % blackboard math symbols
\usepackage{nicefrac}       % compact symbols for 1/2, etc.
\usepackage{microtype}      % microtypography
\usepackage{lipsum}		% Can be removed after putting your text content
\usepackage{graphicx}
\usepackage{natbib}
\usepackage{doi}

\usepackage{amsmath}
\usepackage{amssymb}
\usepackage{mathtools}
\usepackage{amsthm}
\usepackage{multirow}

\usepackage{algorithm}
%\usepackage{algpseudocode}
\usepackage{algorithmic}

%%%%%%%%%%%%%%%%%%%%%%%%%%%%%%%%
% THEOREMS
%%%%%%%%%%%%%%%%%%%%%%%%%%%%%%%%
\theoremstyle{plain}
\newtheorem{theorem}{Theorem}[section]
\newtheorem{proposition}[theorem]{Proposition}
\newtheorem{lemma}[theorem]{Lemma}
\newtheorem{corollary}[theorem]{Corollary}
\theoremstyle{definition}
\newtheorem{definition}[theorem]{Definition}
\newtheorem{example}[theorem]{Example}
\newtheorem{assumption}[theorem]{Assumption}
\theoremstyle{remark}
\newtheorem{remark}[theorem]{Remark}
\usepackage{thm-restate}

\makeatletter
\def\thmt@innercounters{equation,theorem}
\makeatother

\usepackage[colorinlistoftodos]{todonotes}

\newcommand{\sS}{\mathcal{S}}
\newcommand{\sA}{\mathcal{A}}
\newcommand{\sR}{\mathbb{R}}
\newcommand{\train}{\mathrm{train}}
\newcommand{\val}{\mathrm{val}}
\newcommand{\p}{\rho}  % so we can switch the symbol for the stationary distribution
\newcommand{\dist}[1]{\Delta({#1})}
\newcommand{\cD}{\mathcal{D}}
\newcommand{\C}{\mathrm{C}}
\newcommand{\K}{\mathrm{K}}

\definecolor{cbsGreen}{HTML}{00A89D}

\usepackage{tikz}
\usetikzlibrary{bayesnet} % Recommended for Bayesian networks

\usepackage{wrapfig}

\title{Improving State Generalization in Offline Goal-Conditional Reinforcement Learning}

%\date{September 9, 1985}	% Here you can change the date presented in the paper title
%\date{} 					% Or removing it

\author{Nikola Milosevic \\
Max Planck Institute for Human Cognitive and Brain Sciences, Leipzig\\
Center for Scalable Data Analytics and Artificial Intelligence (ScaDS.AI), Dresden/Leipzig\\
\texttt{nmilosevic@cbs.mpg.de} \\
\And
David Carreto Fidalgo \\
Max Planck Computing and Data Facility, Garching\\
\texttt{david.carreto.fidalgo@mpcdf.mpg.de} \\
\And
Nico Scherf \\
Max Planck Institute for Human Cognitive and Brain Sciences, Leipzig\\
Center for Scalable Data Analytics and Artificial Intelligence (ScaDS.AI), Dresden/Leipzig\\
\texttt{nscherf@cbs.mpg.de} \\
}

% Uncomment to remove the date
%\date{}

% Uncomment to override  the `A preprint' in the header
%\renewcommand{\headeright}{Technical Report}
%\renewcommand{\undertitle}{Technical Report}
\renewcommand{\shorttitle}{Geometric Policy Complexity}

%%% Add PDF metadata to help others organize their library
%%% Once the PDF is generated, you can check the metadata with
%%% $ pdfinfo template.pdf
\hypersetup{
pdftitle={A template for the arxiv style},
pdfsubject={q-bio.NC, q-bio.QM},
pdfauthor={Nikola Milosevic, Nico Scherf},
pdfkeywords={Representation Learning, Reinforcement Learning, Offline Reinforcement Learning, Generalization},
}
\def\checkmark{\tikz\fill[scale=0.3](0,.35) -- (.25,0) -- (1,.7) -- (.25,.15) -- cycle;} 

\begin{document}
\maketitle

\begin{abstract}
Representation learning is central to deep reinforcement learning (RL), particularly in high-dimensional, continuous environments, where offline pretraining has proven valuable in improving sample efficiency of model-free algorithms...
\end{abstract}

% keywords can be removed
\keywords{Representation Learning \and Reinforcement Learning \and Offline Reinforcement Learning \and Goal-Conditional Reinforcement Learning}

\section{Introduction} 
Offline Goal-conditional Reinforcement Learning (GCRL) is a popular approach to learning behavior directly from data, without having to interact with an environment for millions of time steps, and without having to design complex reward functions for a variety of different skills. 
While many methods have been proposed (which mostly introduce novel value and policy losses), they all seem to suffer from the same problem: it is difficult to get the policy to generalize to out-of-distribution state visitations at test time. 
When we transfer the learned policy from the pretraining dataset to the online environment, the data we receive from the agent-environment interaction is subject to severe distribution shift compared to the offline dataset, since the training data has been collected with unknown exploration policies. 
It is not clear how to remedy this and when it works, it seems to work mostly because of the "magical" generalization capabilities of deep neural network~\cite{park2024value}. Also, it doesn't really help that most of these methods depend on deep q-learning, which is poorly understood in and of itself. 
We'd like to better understand the mechanisms behind this generalization by performing targeted experiments and trying out different representation learning strategies on top of the offline GCRL losses, since representation learning seems to be a universal tool in helping neural networks generalize. 

\paragraph{Problem.} Assume that we are given a dataset of $N$ trajectories $\mathcal{D} = \{\tau_i\}_{i\in\{1,...,N\}}$ coming from unknown policies executed in an environment. We make no assumptions about the policies or the quality of the data w.r.t a particular task. 
But we are allowed to learn/train anything using this data, but with no access to the environment. In Offline GCRL, no rewards are given in the dataset, but it is known that the task at test time will be a goal-reaching task, hence we can focus on learning goal reaching policies in the pretraining phase.
After pretraining purely on the data, the method must somehow extract a policy for the testing phase, where it is tested on a goal-reaching task or multiple tasks in the online environment. 

\paragraph{Data and Evaluation.} It will make sense to split the trajectory dataset into training and validation sets $\mathcal{D} = (\mathcal{D}_\textrm{train}, \mathcal{D}_\textrm{val})$. 
So far, everything looks a lot like unsupervised machine learning, e.g. we could use plain unsupervised representation learning  techniques to learn state (-action) representations, or learn a generative model over trajectories. 
However, testing is different in our situation: The test dataset is generated dynamically at test time \textit{using the policy we have learned in pretraining}. 
This evaluation protocol introduces a natural distribution shift, even if the environment stays the same: our pretrained policy is different from the ones used to generate the dataset trajectories. Hence, offline RL naturally involves out-of-distribution (OOD) generalization as a sub-problem.

\subsection{Controlled Markov Processes}
We consider the controlled Markov process (CMP) $\mathcal{M} = (\sS, \sA, \mu, p)$, where $\sS$ is the state space, $\sA$ is the action space, $\mu\in \dist\sS$ is the initial state distribution, and $p: \sS\times\sA \to \dist\sS$ is the transition kernel. 
A Markov decision process is obtained from the CMP by adding a reward function $r:\sS\times\sA \to \sR$ to the tuple, i.e. $\mathcal{M}_r = (\sS, \sA, \mu, p, r)$. 
Further, goal-reaching tasks can be expressed as an MDP with a sparse goal-conditioned reward function of the form $r_g(s,a) = \mathbf{1}(\{s=g\})$, where $g\in\mathcal{G}$ is the goal state. 
The goal space $\mathcal{G}$ can be any subset of the state space, but is usually taken to be the entire state space.
An agent interacts with a CMP by choosing a (typically stationary) behavior policy $\pi:\sS\to\dist\sA$. 
Trajectory samples are taken with respect to the initial distribution $s_0\sim\mu$, the policy $a_{t}\sim\pi(\cdot|s_t)$ and the state transition $s_{t+1} \sim p(\cdot|s_t, a_t)$. 
Throughout this work, we will consider collections of distinct policies $\Pi_z = \{\pi_z: z\in \mathcal{Z}\}\subset\Pi$. 
Such a collection will be referred to as a \emph{skill space}.

Given some policy $\pi$ and a reward function $r$, the value function $V_r^\pi\colon \sS \to \mathbb{R}$, action-value function $Q_r^\pi\colon \sS\times\sA \to \mathbb{R}$, and advantage function $A_r^\pi\colon \sS \times \sA \to \mathbb{R}$ are defined as
\begin{equation*}
    V_r^\pi(s) \coloneqq (1-\gamma)\,\mathbb{E}_\pi \left[ \sum_{t=0}^{\infty} \gamma^t r(s_t, a_t) \Big|  s_0 = s \right],    
\end{equation*}
\begin{equation*}
    Q_r^\pi(s, a) \coloneqq (1-\gamma)\,\mathbb{E}_\pi \left[ \sum_{t=0}^{\infty} \gamma^t r(s_t, a_t) \Big| s_0 = s, a_0 = a \right],
\end{equation*}
and
\begin{equation*}
A_r^\pi(s,a) \coloneqq Q_r^\pi(s, a) - V_r^\pi(s).
\end{equation*}
The expectations are taken over trajectories of the Markov process, meaning with respect to the initial distribution $s_0\sim\mu$, the policy $a_{t}\sim\pi(\cdot|s_t)$ and the state transition $s_{t+1} \sim p(\cdot|s_t, a_t)$. Finally $R(\pi) = \mathbb{E}_\mu[V_r(s)]$ is the reward of the policy. All above quantities are defined analogously for any (multivariate) signal $\phi$ instead of $r$, and are denoted as $V_{\phi}^\pi(s)$, $Q_{\phi}^\pi(s,a), A_{\phi}^\pi(s,a)$, and $\psi(\pi)$, respectively, where $\phi$ is often referred to as the cumulant (signal) or (state-action-)features in the literature. Further $V_{\phi}^\pi$, $Q_{\phi}^\pi$ are called successor features, and $\psi(\pi)$ the expected features.

\paragraph{Offline RL} In Offline RL, we are given trajectory data $\mathcal{D} = \{(s_0,a_0,r_0, ...s_T)_i\}_{i\in\{1,...,N\}}$ (including rewards) collected with an arbitrary behavior policy $\pi_b$ in a fixed MDP. The goal is to learn a policy $\pi$ from the data \textit{without environment interaction}, which can be readily executed in the environment without further learning. Possible solution methods are discussed below.

\paragraph{Offline Goal-conditional RL} In Offline GCRL, we are given trajectory data $\mathcal{D} = \{(s_0,a_0, ...s_T)_i\}_{i\in\{1,...,N\}}$ collected with an arbitrary behavior policy $\pi_b$ in a fixed CMP (no rewards!). This time, we are tasked to train a \textit{goal-conditioned} policy $\pi:\sS\times\mathcal{G} \to \dist A$, again, entirely from data. To accomplish this we have to use goal labeling techniques to turn the arbitrary states visited in the trajectory data into imagined goals and fill in rewards accordingly. Sincne most methods involve a q-learning loss, that requires data of just a single time-step, we can randomly choose new goals for each new state transition sample. The goal is usually sampled from the same trajectory that the current state transition is from. Given a trajectory $\tau$, we sample a state-action-next-state tuple $(s,a,s')\sim\tau$, and a goal state $g\sim\tau$ from the same trajectory according to one of the following distributions:
\begin{itemize}
    \item $p_\delta(g|s) = \mathbf{1}(g=s)$
    \item $p_\mathrm{future}(g|s) = \mathbf{1}(g=s_k)$, $k\sim\mathrm{Unif}(\min(t+1,T-1), T-1)$
    \item $p_\mathrm{geom}(g|s) = \mathbf{1}(g=s_{\min(s+k, T-1)})$, $k\sim\mathrm{Geom}(1-\gamma)$
\end{itemize}

\paragraph{Challenges} Learned policy can deviate too much from the policies used to generate the data, which will result in entirely different state visitations online. The ability to learn (a generalizable) value function can be limited due to function approximation error. Policy extraction from the learned value function can be limited again due to function approximation error or a sub-optimally chose loss. Policy needs to generalize to new state online.

\section{Improving State Generalization in Offline GCRL}

We will operate under the working hypothesis that state representation is the "working horse" behind these models: Ultimately, our neural network policy is trained to predict actions from states, and regardless of what losses we used to train it, it's features must support proper state generalization for predicting the actions correctly in the online test, including generalization between (abstract) goal states $g$ or $z$.

\paragraph{Empirical Idea: Representation Visualization} Make visualizations to track ReLU activation patterns during trajectories in the dataset vs. trajectories on the test environment. Compare value functions vs. policy networks. Start with the direct state environments (not visual). Can we understand how the state space is partitioned into linear control regions, and what are the controller gains there?

\paragraph{Improvement Idea} The policy must be trained conservatively, but not with respect to the states in the dataset but w.r.t the states the new policy is expected to visit in the long run.


\section{Experiments}
\label{sec:experiments}

\bibliographystyle{unsrtnat}
\bibliography{references}

\appendix

\section{Experiment Details}
\subsection{Implementation Details}
\subsection{Ogbench Environmets}
\paragraph{Regular Datasets}
\begin{itemize}
    \item antmaze-{large, giant}-navigate
    \item humanoidmaze-medium-navigate
    \item cube-{single, double}-play
    \item scene-play
    \item puzzle-3x3-play
    \item \textbf{Good Algorithms:} GCIVL/GCIQL for Manipulation with states/pixels and for Drawing, CRL for locomotion
\end{itemize}

\paragraph{Harder Tasks}
\begin{itemize}
    \item humanoidmaze-{large, giant}
    \item antsoccer
    \item puzzle-{4x4, 4x5, 4x6}
    \item powderworld
    \item \textbf{Good Algorithms:} GCIVL/GCIQL for Manipulation with states/pixels and for Drawing, CRL for locomotion
\end{itemize}

We also provide more specialized datasets that pose specific challenges in offline GCRL.

\paragraph{Stitching Tasks}
\begin{itemize}
    \item humanoidmaze-{large, giant}-stitch
    \item puzzle-{4x4, 4x5, 4x6}
    \item \textbf{Good Algorithms:} from states all are good, from pixels CRL is good
\end{itemize}

\paragraph{Long-Horizon Tasks}
\begin{itemize}
    \item humanoidmaze-giant
    \item puzzle-4x6
    \item \textbf{Good Algorithm:} GCIVL
\end{itemize}

\paragraph{Stochastic}
\begin{itemize}
    \item powderworld
    \item antmaze-teleport
    \item \textbf{Good Algorithm:} CRL
\end{itemize}

\paragraph{Exploratory (Highly Suboptimal) Data}
\begin{itemize}
    \item antmaze-explore
    \item {puzzle, scene, cube}-noisy
    \item \textbf{Good Algorithm:} GCIVL
\end{itemize}
For learning from highly suboptimal data, consider antmaze-explore as well as the noisy datasets in the manipulation suite, which feature high suboptimality and high coverage.

\end{document}
